\section{Annotation Pipeline}
\label{sec:annotation}

\subsection{Annotation Platform}
\label{ssec:ann_platform}

We developed a bespoke web-based annotation tool, \textit{LUCID-Annotate}, designed
specifically for causal and counterfactual annotation of driving scenarios. The tool
renders synchronized multi-camera video alongside LiDAR depth projections, trajectory
overlays, and the pre-annotated 3D object tracks from the preprocessing pipeline. Key
features of LUCID-Annotate include:

\begin{itemize}[leftmargin=2em, noitemsep]
  \item \textbf{Causal graph editor}: An interactive directed-graph canvas where annotators
        create event nodes by selecting time windows and describing observable events,
        then draw directed edges to represent causal relationships.
  \item \textbf{Counterfactual wizard}: A guided interface that presents an existing causal
        QA pair and prompts the annotator to specify a counterfactual premise and answer.
        Consistency checks warn when the entered premise violates physical constraints.
  \item \textbf{Agent intention panel}: A structured form listing all tracked agents with
        dropdown menus for intention labels and a free-text field for ambiguous cases.
  \item \textbf{Risk matrix tool}: A structured risk-assessment form with numerical sliders
        (hazard severity 1--5, probability 0--1) linked to a free-text explanation field.
  \item \textbf{Review interface}: A dedicated quality-control view presenting annotations
        from two annotators side by side for adjudication.
\end{itemize}

\subsection{Annotator Recruitment and Training}
\label{ssec:ann_recruitment}

A total of 234 annotators were recruited across three annotation centres (Shanghai,
Munich, and London) to ensure geographic representation in annotation decisions.
Annotator selection required: (a) native or near-native English proficiency; (b) a
valid driving licence with at least three years of active driving experience; and
(c) successful completion of a domain-knowledge assessment covering traffic rules,
common hazard scenarios, and basic causal reasoning concepts.

Training consisted of a structured five-day programme:
\begin{enumerate}[leftmargin=2em, noitemsep]
  \item \textbf{Day 1}: Platform orientation and annotation schema walkthrough.
  \item \textbf{Day 2}: Causal chain annotation practice on 20 gold-standard scenarios.
  \item \textbf{Day 3}: Counterfactual reasoning training with explicit examples of
        valid and invalid counterfactual hypotheses.
  \item \textbf{Day 4}: Risk assessment and planning instruction annotation practice.
  \item \textbf{Day 5}: Calibration session --- annotators annotate 10 scenarios independently
        and results are reviewed with supervisors. Annotators with agreement score
        $\kappa < 0.65$ with the gold standard received supplementary training.
\end{enumerate}

\subsection{Annotation Workflow}
\label{ssec:ann_workflow}

Each scenario was assigned to three independent annotators following a structured workflow:

\begin{enumerate}[leftmargin=2em, noitemsep]
  \item \textbf{Stage 1 --- Primary Annotation}: Two primary annotators independently
        complete all five annotation levels for their assigned scenario. Estimated time
        per scenario: 22.4~minutes (mean).
  \item \textbf{Stage 2 --- Cross-review}: A third annotator reviews both primary annotations,
        identifies discrepancies, and proposes a merged resolution using majority voting
        for structured fields and paraphrase alignment for free-text fields.
  \item \textbf{Stage 3 --- Expert Validation}: For causal chain annotations with disagreement
        rate $> 25\%$ and for all counterfactual pairs, a domain expert (senior engineer
        with $> 5$ years in AD) performs a final validation pass. Expert validators
        reviewed 18.7\% of all scenarios.
\end{enumerate}

\subsection{Quality Control Metrics}
\label{ssec:ann_quality}

We computed inter-annotator agreement (IAA) at each annotation level:

\begin{itemize}[leftmargin=2em, noitemsep]
  \item \textbf{Agent intention labels}: Cohen's $\kappa = 0.84$, indicating strong
        agreement across the 27-class intention vocabulary.
  \item \textbf{Causal chain structure}: Structural similarity of causal DAGs measured
        via graph edit distance normalised by graph size; mean similarity = $0.81$.
  \item \textbf{Causal QA pair answers}: BERTScore F1 between independently authored
        causal answers = $0.87$ (threshold for acceptance: $> 0.75$).
  \item \textbf{Counterfactual QA answers}: BERTScore F1 = $0.79$.
  \item \textbf{Risk assessment scores}: Intraclass correlation coefficient (ICC) = $0.88$
        for hazard severity; ICC = $0.82$ for probability.
\end{itemize}

The overall annotation campaign required approximately 45,000 person-hours over
14 months (the first 4 months were dedicated to platform development and annotator training).

\subsection{Annotation Statistics}
\label{ssec:ann_stats}

Table~\ref{tab:annotation_stats} summarises the annotation production statistics.

\begin{table}[htbp]
  \centering
  \caption{Summary of \lucid{} annotation production statistics.}
  \label{tab:annotation_stats}
  \small
  \begin{tabular}{lrrr}
    \toprule
    \textbf{Annotation Type} & \textbf{Count} & \textbf{Mean per Scenario} & \textbf{Avg.\ Length (words)} \\
    \midrule
    Scene descriptions           & 89,247  & 6.9  & 34.2 \\
    Agent intention labels       & 412,847 & 32.1 & --   \\
    Causal QA pairs              & 312,456 & 24.3 & 58.7 \\
    Counterfactual QA pairs      & 198,327 & 15.4 & 52.3 \\
    Planning instruction pairs   & 156,832 & 12.2 & 47.1 \\
    Risk assessment annotations  & 90,431  & 7.0  & 41.8 \\
    \midrule
    \textbf{Total language annot.}& \textbf{847,293} & \textbf{65.9} & \textbf{50.1} \\
    \bottomrule
  \end{tabular}
\end{table}
